\section{Conclusion}

We studied Replayed-Prefix On-Policy Distillation in an efficient, off-environment regime
that reuses a fixed pool of pre-collected teacher trajectories, replaying a
teacher-forced prefix and supervising the student's action at the current step.
Removing the environment does not remove the temporal structure of multi-turn learning:
we identified a \emph{prefix trap} with a temporal layer (compounding errors) and a
distributional layer (a two-sided shift), and bounded the gap to an ideal interactive
improvement objective~(\plaineqref{eq:ideal}) by a student occupancy-mismatch term and a
teacher reliability term~(\plaineqref{eq:bound}). This recasts multi-turn OPD as \emph{reliability-aware prefix distribution design}
-- balancing student relevance against teacher reliability, with fully
student-on-policy and teacher-forced roll-ins as the two extremes -- which the
resulting algorithm, ReOPD, realizes with a simple step-decay schedule applied by
default as a sampling probability (equivalently, a loss weight). Our experiments confirm the predicted
regime-aware behavior: ReOPD stays OPD-like when the teacher--student gap is small and
yields its largest gains when the gap is large and teacher-supported prefixes are the
more reliable supervision region.

\paragraph{Limitations and future work.}
Our analysis assumes access to a pre-collected pool of teacher trajectories with
recorded observations. In our experiments this comes for free from the teacher's
own RL rollouts, but more generally the pool's coverage and quality bound what the
student can learn. The reliability surrogate we adopt reduces, to first order, to a
step-decaying weight, which is effective but coarse: it captures \emph{how deep} a
prefix is rather than directly measuring \emph{how far} a given prefix lies from the
overlap between student-relevant histories and the teacher's reliable support, and a
learned or data-dependent reliability estimate may improve on the schedule. Finally,
our two-sided bound treats the teacher reliability error through teacher support as a
surrogate; tighter, directly estimable notions of teacher reliability -- and schedules
or weights that adapt to them online -- are a natural next step.
